% !TeX root = ../main.tex
\documentclass[
    ngerman,
    accentcolor=3b,
    dark_mode,
    fontsize= 12pt,
    a4paper,
    aspectratio=169,
    colorback=true,
    fancy_row_colors,
    leqno,
    fleqn,
    boxarc=3pt,
    fleqn,
    % shell_escape = false, % Kompatibilität mit sharelatex
]{algoslides}

%%--------------------------%%
%%--Imports from Main File--%%
%%--------------------------%%
\usepackage{import}
% Import all Packages from Main Preamble with relative Path (buggy, list packages instead)
\subimport*{../shared}{preamble}
% Get Labels from Main Document using the xr-hyper Package
\externaldocument[ext:]{../main}
% Set Graphics Path, so pictures load correctly
\graphicspath{{../pictures}}
\def\codeDir{../code}

\begin{document}
    \section{Abbildungen}
    \subsection{Bilder einfügen}
    \begin{frame}[c, fragile]
        \slidehead{}
        %\begin{noindent}
        \begin{codeBlock}[
            % highlightlines={7}
            ]{
            minted language=latex,
            listing and text,
            sidebyside,
            righthand width=2.1cm,
            %center lower,
            bicolor,
            colbacklower=\IfDarkModeTF{codebg!90}{codebg!20},
            title=\codeBlockTitle{Einfügen einer Abbildung},
            }
            \begin{figure}[ht!]
                \centering
                \includegraphics[width=2cm, height=2cm, keepaspectratio]{LaTeX-Wallpaper.png}
                \caption{LaTeX Wallpaper}
                \label{fig:latexlogo}
            \end{figure}
        \end{codeBlock}
        %\end{noindent}
        \begin{itemize}
            \item References in Kapitel 7 erklärt
        \end{itemize}
    \end{frame}
    \begin{frame}[c, fragile]
        \slidehead{}
        \vspace{-2em}
        \footnotesize
        %\begin{noindent}
        \begin{codeBlock}[
            % highlightlines={7}
            ]{
            minted language=latex,
            listing and text,
            sidebyside,
            righthand width=4.1cm,
            %center lower,
            bicolor,
            colbacklower=\IfDarkModeTF{codebg!90}{codebg!20},
            title=\codeBlockTitle{Subfigures},
            top=0pt,
            bottom=0pt,
            }
            \begin{figure}[ht!]
                \centering
                \begin{subfigure}[c]{.5\textwidth}
                    \centering
            \includegraphics[width=2cm]{LaTeX-Wallpaper.png}
                    \caption{LaTeX Wallpaper}
                    \label{fig:pictures:latexwp}
                \end{subfigure}%
                \begin{subfigure}[c]{.5\textwidth}
                    \centering
                    \includegraphics[width=2cm]{dsicord.png}
                    \caption{Discord Logo}
                    \label{fig:pictures:discord}
                \end{subfigure}
                \caption{Bilder}
                \label{fig:pictures}
            \end{figure}
        \end{codeBlock}
        %\end{noindent}
        \begin{itemize}
            \item References in Kapitel 7 erklärt
        \end{itemize}
    \end{frame}
    \begin{frame}[c, fragile]
        \slidehead{}
        \vspace{-2em}
        \footnotesize
        %\begin{noindent}
        \begin{codeBlock}[
            % highlightlines={7}
            ]{
            minted language=latex,
            listing and text,
            sidebyside,
            righthand width=4.1cm,
            %center lower,
            bicolor,
            colbacklower=\IfDarkModeTF{codebg!90}{codebg!20},
            title=\codeBlockTitle{Subfigures},
            top=0pt,
            bottom=0pt,
            }
            \begin{figure}[ht!]
                \centering
                \begin{subfigure}[b]{.5\textwidth}
                    \centering
            \includegraphics[width=2cm]{LaTeX-Wallpaper.png}
                    \caption{LaTeX Wallpaper}
                    \label{fig:pictures:latexwp}
                \end{subfigure}%
                \begin{subfigure}[b]{.5\textwidth}
                    \centering
                    \includegraphics[width=2cm]{dsicord.png}
                    \caption{Discord Logo}
                    \label{fig:pictures:discord}
                \end{subfigure}
                \caption{Bilder}
                \label{fig:pictures}
            \end{figure}
        \end{codeBlock}
        %\end{noindent}
        \begin{itemize}
            \item References in Kapitel 7 erklärt
        \end{itemize}
    \end{frame}
    \begin{frame}[c, fragile]
        \slidehead{}
        \vspace{-2em}
        \footnotesize
        %\begin{noindent}
        \begin{codeBlock}[
            % highlightlines={7}
            ]{
            minted language=latex,
            listing and text,
            sidebyside,
            righthand width=4.1cm,
            %center lower,
            bicolor,
            colbacklower=\IfDarkModeTF{codebg!90}{codebg!20},
            title=\codeBlockTitle{Subfigures},
            top=0pt,
            bottom=0pt,
            }
            \begin{figure}[ht!]
                \centering
                \begin{subfigure}[t]{.5\textwidth}
                    \centering
            \includegraphics[width=2cm]{LaTeX-Wallpaper.png}
                    \caption{LaTeX Wallpaper}
                    \label{fig:pictures:latexwp}
                \end{subfigure}%
                \begin{subfigure}[t]{.5\textwidth}
                    \centering
                    \includegraphics[width=2cm]{dsicord.png}
                    \caption{Discord Logo}
                    \label{fig:pictures:discord}
                \end{subfigure}
                \caption{Bilder}
                \label{fig:pictures}
            \end{figure}
        \end{codeBlock}
        %\end{noindent}
        \begin{itemize}
            \item References in Kapitel 7 erklärt
        \end{itemize}
    \end{frame}
    \begin{frame}[c, fragile]
        \slidehead{}
        \vspace{-2em}
        \footnotesize
        %\begin{noindent}
        \begin{codeBlock}[
            % highlightlines={7}
            ]{
            minted language=latex,
            listing and text,
            sidebyside,
            righthand width=4.1cm,
            %center lower,
            bicolor,
            colbacklower=\IfDarkModeTF{codebg!90}{codebg!20},
            title=\codeBlockTitle{Subfigures},
            top=0pt,
            bottom=0pt,
            }
            \begin{figure}[ht!]
                \centering
                \begin{subfigure}[t]{.5\textwidth}
                    \centering\mbox{}
            \includegraphics[width=2cm]{LaTeX-Wallpaper.png}
                    \caption{LaTeX Wallpaper}
                    \label{fig:pictures:latexwp}
                \end{subfigure}%
                \begin{subfigure}[t]{.5\textwidth}
                    \centering\mbox{}
                    \includegraphics[width=2cm]{dsicord.png}
                    \caption{Discord Logo}
                    \label{fig:pictures:discord}
                \end{subfigure}
                \caption{Bilder}
                \label{fig:pictures}
            \end{figure}
        \end{codeBlock}
        %\end{noindent}
        \begin{itemize}
            \item References in Kapitel 7 erklärt
        \end{itemize}
    \end{frame}
    \subsection{Tabellen}
    \begin{frame}[c, fragile]
        \slidehead{}
        \vspace{-2em}
        \footnotesize
        %\begin{noindent}
        \begin{codeBlock}[
            % highlightlines={7}
            ]{
            minted language=latex,
            listing and text,
            sidebyside,
            %righthand width=4.1cm,
            %center lower,
            bicolor,
            colbacklower=\IfDarkModeTF{codebg!90}{codebg!20},
            title=\codeBlockTitle{Subfigures},
            top=0pt,
            bottom=0pt,
            }
            \begin{table}
                \begin{tabular}{lcrp{2cm}}
                    \toprule
                    \textbf{Spalte 1} & \textbf{Spalte 2} & \textbf{Spalte 3} & \textbf{Spalte 4} \\
                    \midrule
                    1 & 2 & 3 & 4 \\
                    5 & 6 & 7 & 8 \\
                    9 & 10 & 11 & 12 \\
                    \bottomrule
                \end{tabular}
                \caption{Eine Tabelle}
                \label{tab:table}
            \end{table}
        \end{codeBlock}
        %\end{noindent}
    \end{frame}
    % \subsection{TikZ}
\end{document}
